\chapter{Architecture}
\label{ch:arch}

\section{Module map}
\begin{center}
\begin{tikzpicture}[node distance=7mm and 11mm,font=\footnotesize]
  \node[fnblockD,minimum width=34mm,minimum height=13mm] (dm){\textbf{Dependency Manager}\\{\scriptsize node table, wake-up}};
  \node[fnblockT,right=13mm of dm,minimum width=32mm,minimum height=13mm] (dir){\textbf{Neural Director}\\{\scriptsize first-free dispatch}};
  \node[fnreg,fill=white,right=13mm of dir,minimum width=34mm,minimum height=20mm] (mm0){
    \begin{tabular}{c}\textbf{Memory Manager} 0\\ $+$ \textbf{Neural Processor} 0\end{tabular}};
  \node[fnreg,fill=white,below=3mm of mm0,minimum width=34mm,minimum height=20mm] (mm1){
    \begin{tabular}{c}\textbf{Memory Manager} 1\\ $+$ \textbf{Neural Processor} 1\end{tabular}};
  \node[fnblockA,below=9mm of dir,minimum width=32mm,minimum height=13mm] (cache){\textbf{Activation Cache}\\{\scriptsize shared, single-tag}};
  \node[fnblock,right=13mm of mm0,minimum width=26mm,minimum height=13mm] (arb){\textbf{Slot Memory}\\\textbf{Arbiter}};
  \node[fnblockD,below right=9mm and 13mm of arb,minimum width=30mm,minimum height=13mm] (psram){\textbf{Real V1 PSRAM chain}\\{\scriptsize memory\_interface $\to$ psram\_controller}};
  \draw[fnbus] (dm) -- node[fnlbl,above]{ready\_valid/ready} (dir);
  \draw[fnbus] (dir) -- (mm0);
  \draw[fnbus] (dir) -- (mm1);
  \draw[fnarrowT] (mm0.south) |- (cache.east);
  \draw[fnarrowT] (mm1.west) -- (cache.east);
  \draw[fnarrowT] (cache.north) |- node[fnlbl,above,pos=0.3]{producer\_done} (dm.south);
  \draw[fnbus] (mm0) -- (arb);
  \draw[fnbus] (mm1) -- (arb);
  \draw[fnbus] (cache.south) |- (arb.west);
  \draw[fnbus] (arb) -- (psram);
\end{tikzpicture}
\end{center}
\begin{center}\footnotesize\itshape\color{fnGrey}
N\_SLOTS=2 shown (the recommended configuration); the architecture is
parametric in N\_SLOTS. Every arrow is a real signal path verified in
Verilator simulation and real Yosys/nextpnr-ecp5 synthesis.\end{center}

\section{Dependency Manager}
Holds a table of \code{N\_NODES} job descriptors, each tracking: node
id, state (\code{EMPTY}/\code{WAITING}/\code{READY}/\code{DISPATCHED}),
required-dependency count, resolved-dependency count, up to
\code{MAX\_DEPS} producer node ids, and the job descriptor fields
(\code{x\_base}, \code{w\_base}, \code{n\_tiles}, \code{result\_addr}). A
node with zero required dependencies is immediately \code{READY} on
registration. When a producer completes, \emph{every} \code{WAITING}
node listing it among its own producers gets its resolved-dependency
count incremented --- a single producer can satisfy several waiting
consumers (shared-producer/multi-consumer), and a node depending on
several producers accumulates resolution across separate events
(multiple dependencies). Verified for both 1-hop and 2-hop transitive
(diamond) graphs. Ready nodes are handed to the Neural Director one at a
time over a backpressure-safe valid/ready interface.

\begin{fnwarn}[No slot reclamation]
\code{ST\_DISPATCHED} is terminal: node table slots are never reused
once dispatched. A long-running system that keeps registering new
nodes without limit will eventually exhaust \code{N\_NODES} --- this is
a real, measured consequence (a benchmark testbench hit exactly this
deadlock via node-id wraparound before \code{N\_NODES} was sized
generously enough). Slot reclamation is explicitly deferred, not
forgotten.
\end{fnwarn}

\section{Neural Director}
Dispatches ready job descriptors to whichever of \code{N\_SLOTS} Memory
Manager instances is currently free --- \textbf{first-free} scheduling: a
fixed, lowest-index-wins priority scan, not load-balanced. Slot-busy
tracking and completion detection are always-active, independent of
whatever the allocate/scan control state happens to be that cycle (the
same ``don't gate a per-unit event behind one shared FSM state''
principle applied throughout this design). A completed slot's node id
is tracked (\code{slot\_node\_id}) so its completion can be resolved back
to a \code{producer\_done} event for the Dependency Manager, closing the
wake-up loop without any external glue logic.

\begin{fnnote}[Measured scheduling imbalance]
Real per-slot data (\code{N\_SLOTS}=4, a 128-neuron workload) shows
slots 0/1 delivering 1008 real tiles each while slots 2/3 deliver only
16 each, despite all four slots reporting near-100\% ``busy''
utilization --- direct, measured evidence that fixed lowest-index
priority does not distribute load evenly once the shared PSRAM port is
the real constraint. See ch.~\ref{ch:impl2}.
\end{fnnote}

\section{Memory Manager + Neural Processor (per slot)}
Each slot pairs one \code{memory\_manager.v} instance with one
\code{neural\_processor.v} instance. The Memory Manager double-buffers
tile fetches (compute tile $N$ while prefetching tile $N{+}1$) and
presents the Neural Processor with a simple ``data available''
interface (\code{operand\_valid/ready}, \code{tile\_last}) --- the
processor never sees PSRAM request/wait cycles directly. A tile's
activation half is requested from the shared Activation Cache; its
weight half is fetched directly (weights are per-neuron, never shared,
so caching them would not help). A bank is presentable to the processor
only once \emph{both} halves have arrived
(\code{bank\_ready = bank\_x\_ready \& bank\_w\_ready}).

\section{Activation Cache}
\label{sec:archcache}
A single shared instance (not one per slot) serving every Memory
Manager's activation-fetch requests. Single-tag design: one cached
\code{x\_base} at a time, filled tile-by-tile on first use, served
directly from an on-chip buffer on every subsequent request for the same
vector --- no PSRAM access on a hit. A request for a different
\code{x\_base} invalidates the cache and restarts filling from tile~0;
this is always \emph{correct} (never serves stale data) but can thrash
under interleaved, genuinely-different-\code{x\_base} concurrent
traffic --- an honestly documented limitation, not exercised by this
project's own realistic dense-layer workloads (where many neurons of one
layer share one input vector, dispatched together). Full detail,
including the real Fmax cost this module introduces, in
ch.~\ref{ch:mem}.

\section{Slot Memory Arbiter}
Funnels \code{N\_SLOTS}$+$1 independent backend ports (one per Memory
Manager's weight/write-back traffic, plus one for the Activation
Cache's own traffic) down to the one real physical PSRAM port. Fixed
lowest-index priority, same convention as the Director. Every incoming
request is latched into a per-port pending register regardless of
arbiter state --- a byte-level backend protocol quirk discovered by real
simulation (a fire-and-forget single-cycle request pulse can arrive
while the shared bus is owned by another port; a naive ``grant only
while live'' arbiter would silently drop it) made this latch a
correctness requirement, not an optimization.

\section{Real, unmodified V1 PSRAM backend}
\code{memory\_interface.v} and \code{psram\_controller.v} are reused
byte-for-byte from the frozen \code{hardware/v1/} tree. The controller's
own real page-mode support (fast same-page continuation vs.\ a slower
cold access) was already implemented in V1 and is exploited more
effectively by V2's word-level burst rewrite (ch.~\ref{ch:mem}) --- no
change to the controller itself was needed or made.
